\documentclass[12pt,a4paper]{article}
\usepackage[utf8]{inputenc}
\usepackage[T1]{fontenc}
\usepackage[lithuanian]{babel}
\usepackage{lmodern}
\usepackage{amsmath, amssymb}
\usepackage{geometry}
\geometry{left=2cm, right=2cm, top=2.5cm, bottom=2.5cm}

\begin{document}

\begin{center}
    {\Large \textbf{Kontrolinis darbas: Daugianariai}}\\[0.3cm]
    {\large \textbf{Aristotelio variantas -- Atsakymai ir sprendimai}}
\end{center}
\vspace{0.5cm}

\begin{enumerate}
    \item \textbf{Reiškinių tipai:}
    \begin{itemize}
        \item[a)] $\displaystyle \frac{3-5x}{4}$ -- \textbf{R, S} (Racionalusis sveikas). Vardiklyje nėra kintamojo, dalijama tik iš skaičiaus.
        \item[b)] $\displaystyle \frac{9}{x^2 + 1}$ -- \textbf{R, T} (Racionalusis trupmeninis). Kintamasis vardiklyje.
        \item[c)] $\displaystyle 3x^2 + \sqrt{2} x$ -- \textbf{R, S} (Racionalusis sveikas). Šaknis traukiama tik iš skaičiaus $2$.
        \item[d)] $\displaystyle \sqrt{x^3 + 1}$ -- \textbf{I} (Iracionalusis). Kintamasis yra po šaknies ženklu.
        \item[e)] $\displaystyle x + 3^{-2}$ -- \textbf{R, S} (Racionalusis sveikas). Skaičius $3^{-2} = \frac{1}{9}$, kintamasis nėra vardiklyje.
        \item[f)] $\displaystyle \sqrt{6}x - 7x^2$ -- \textbf{R, S} (Racionalusis sveikas). Šaknis traukiama tik iš skaičiaus $6$.
    \end{itemize}

    \vspace{0.3cm}
    \item \textbf{Vienanarių daugyba:} \\
    Pritaikome laipsnių savybes:
    $$ \left( -2x^3y^{-4} \right)^3 \cdot \left( 2x^{-6}y^5 \right)^{-2} = (-2)^3 \cdot x^9 \cdot y^{-12} \cdot 2^{-2} \cdot x^{12} \cdot y^{-10} $$
    $$ = -8 \cdot x^9 \cdot y^{-12} \cdot \frac{1}{4} \cdot x^{12} \cdot y^{-10} = -2 \cdot x^{9+12} \cdot y^{-12-10} = \mathbf{-2x^{21}y^{-22}} \text{ arba } \mathbf{-\frac{2x^{21}}{y^{22}}} $$
    \textbf{Paaiškinimas:} Gautas reiškinys \textbf{nėra vienanaris}, nes kintamojo $y$ laipsnio rodiklis yra neigiamas (tai reiškia dalybą iš kintamojo, o vienanariuose kintamųjų dalyti negalima).

    \vspace{0.3cm}
    \item \textbf{Skaičiavimas be skaičiuotuvo:} \\
    \begin{align*}
        \frac{167^3 + 133\left(121^2 + 2 \cdot 12 \cdot 121 + 12^2\right)}
        {167(10\sqrt{3})^2\left(34 + \frac{133^2}{167}\right)}
        &= \frac{167^3 + 133 \cdot 133^2}
        {300 \cdot 167\left(34 + \frac{133^2}{167}\right)} \\
        &= \frac{167^3 + 133^3}
        {300\left(167 \cdot 34 + 133^2\right)} \\
        &= \frac{(167+133)\left(167^2-167 \cdot 133+133^2\right)}
        {300\left(167 \cdot 34 + 133^2\right)} \\
        &= \frac{(167+133)\left(167(167-133)+133^2\right)}
        {300\left(167 \cdot 34 + 133^2\right)} \\
        &= \frac{300\left(167 \cdot 34 + 133^2\right)}
        {300\left(167 \cdot 34 + 133^2\right)} = 1.
    \end{align*}

    \vspace{0.3cm}
    \item \textbf{Formulių įrodymai:}
    \begin{itemize}
        \item[a)] $\displaystyle \frac{x^m}{x^n} = \frac{\overbrace{x \cdot x \cdot \dots \cdot x}^{m \text{ kartų}}}{\underbrace{x \cdot x \cdot \dots \cdot x}_{n \text{ kartų}}} = \underbrace{x \cdot x \cdot \dots \cdot x}_{m-n \text{ kartų}} = \mathbf{x^{m-n}}.$ (kai $m > n$).
        \item[b)] Atskliaudžiame dešinę pusę: \\
        $(x-y)(x^2 + xy + y^2) = x^3 + x^2y + xy^2 - x^2y - xy^2 - y^3 = \mathbf{x^3 - y^3}.$
    \end{itemize}

    \vspace{0.3cm}
    \item \textbf{Didžiausia ir mažiausia reikšmės:} \\
    Išskiriame pilnąjį kvadratą:
    $$ -4x^2 - x + 1 = -4\left(x^2 + \frac{1}{4}x\right) + 1 = -4\left(x^2 + 2 \cdot x \cdot \frac{1}{8} + \frac{1}{64} - \frac{1}{64}\right) + 1 $$
    $$ = -4\left( \left(x + \frac{1}{8}\right)^2 - \frac{1}{64} \right) + 1 = -4\left(x + \frac{1}{8}\right)^2 + \frac{4}{64} + 1 = \mathbf{-4\left(x + \frac{1}{8}\right)^2 + 1\frac{1}{16}} $$
    Kadangi parabolės šakos nukreiptos žemyn, \textbf{didžiausia reikšmė yra $\mathbf{1\frac{1}{16}}$} (arba $\frac{17}{16}$), kai $x = -\frac{1}{8}$. \\
    \textbf{Mažiausios reikšmės nėra} (artėja į $-\infty$).

    \vspace{0.3cm}
    \item \textbf{Daugianarių dalyba ir laipsnis:} \\
    Atskliaudžiame ir suprastiname reiškinį $(2x+1)^3 - (4x-1)(x+1)^2 - 4\left(x+\frac{1}{2}\right)$: \\
    $ (8x^3 + 12x^2 + 6x + 1) - (4x-1)(x^2 + 2x + 1) - (4x + 2) $ \\
    $ = 8x^3 + 12x^2 + 6x + 1 - (4x^3 + 8x^2 + 4x - x^2 - 2x - 1) - 4x - 2 $ \\
    $ = 8x^3 + 12x^2 + 6x + 1 - (4x^3 + 7x^2 + 2x - 1) - 4x - 2 $ \\
    $ = 8x^3 + 12x^2 + 6x + 1 - 4x^3 - 7x^2 - 2x + 1 - 4x - 2 $ \\
    Sutraukiame panašiuosius narius: \\
    $ = (8-4)x^3 + (12-7)x^2 + (6-2-4)x + (1+1-2) = \mathbf{4x^3 + 5x^2} $ \\
    Padaliję šį reiškinį iš $x^2$, gauname:
    $$ \frac{4x^3 + 5x^2}{x^2} = \mathbf{4x + 5} $$
    Tai yra daugianaris, jo \textbf{laipsnis lygus 1}.

    \vspace{0.3cm}
    \item \textbf{Nežinomų koeficientų radimas:} \\
    Atskliaudžiame kairiąją lygybės pusę:
    $$ (2x-3)(3x^2+kx-1) + m = 6x^3 + 2kx^2 - 2x - 9x^2 - 3kx + 3 + m $$
    $$ = 6x^3 + (2k-9)x^2 + (-2-3k)x + (3+m) $$
    Sulyginame su dešine puse ($6x^3 - x^2 - 14x + 10$):
    \begin{itemize}
        \item Prie $x^2$: $2k - 9 = -1 \implies 2k = 8 \implies \mathbf{k = 4}$.
        \item Patikriname prie $x$: $-2 - 3(4) = -14$ (teisinga).
        \item Laisvasis narys: $3 + m = 10 \implies \mathbf{m = 7}$.
    \end{itemize}
    Skaičiuojame prašomo reiškinio reikšmę:
    $$ \left(-\frac{1}{4} \cdot 4\right)^{666} + \left(-\frac{1}{7} \cdot 7\right)^{999} = (-1)^{666} + (-1)^{999} = 1 + (-1) = \mathbf{0} $$

\end{enumerate}

\end{document}